\chapter[Lorentz locking]{Local Lorentz Locking and the Fate of the Goldstones}
\label{ch:lorentz-locking}

\paragraph{Learning outcomes.}
After this chapter the reader should be able to:
\begin{enumerate}
  \item write the simultaneous action of internal gauge
        transformations and local Lorentz transformations on
        \(\Phi^I{}_A\);
  \item identify the diagonal subgroup preserved by the polar vacuum
        and the six broken generators;
  \item exhibit the Stueckelberg replacement that eats
        \(\omega_{AB}\) into a spin connection;
  \item finish the counting table of Section~\ref{sec:ch5-count}
        through the vector and Lorentz-Goldstone lines;
  \item enter Chapter~11 with a composite frame that transforms as a
        genuine vierbein, so that \(\nabla^\mu T_{\mu\nu}\) is well
        defined.
\end{enumerate}

Since Chapter~5 the antisymmetric piece \(\omega_{AB}\) has been
called a local-Lorentz Goldstone and postponed. This chapter
constructs the gauge symmetry that justifies the name and the Higgs
mechanism that removes \(\omega_{AB}\) from the infrared spectrum.

\section{The two groups that act on \(\Phi\)}
\label{sec:ch10-groups}

The condensate field of Presentation~B transforms, by
\eqref{eq:Phi-transform}, as
\begin{equation}
  \label{eq:Phi-tr-again}
  \Phi^I{}_A(x)
  \;\longmapsto\;
  \Omega^I{}_J(x)\,\Phi^J{}_B(x)\,\Lambda^B{}_A(x).
\end{equation}
Here \(\Omega(x)\) takes values in the internal gauge group
\(G_{\mathrm{int}}\) and \(\Lambda(x)\) takes values in
\(\mathrm{SO}(3,1)\), or in its double cover if spinors are later
coupled. Infinitesimally,
\begin{equation}
  \label{eq:Phi-inf}
  \delta\Phi^I{}_A
  =\alpha^I{}_J\Phi^J{}_A
   +\Phi^I{}_B\,\beta^B{}_A,
\end{equation}
with \(\alpha\) in the Lie algebra of \(G_{\mathrm{int}}\) and
\(\beta_{AB}=-\beta_{BA}\) in \(\mathfrak{so}(3,1)\).

These are independent transformations at the level of field content.
Nothing in Chapter~3 forced \(\Lambda(x)\) to be a gauge symmetry
rather than a global classification of the index \(A\). Treating it
as local is a choice, recorded as such.

\begin{assumption}[Local Lorentz gauge symmetry]
  \label{ass:local-lorentz}
  The microscopic action~\eqref{eq:Smicro} is invariant under
  local \(\Lambda(x)\in\mathrm{SO}(3,1)\) once a spin connection
  \(\omega_\mu{}^A{}_B\) is introduced through the covariant
  derivative on the tangent index,
  \begin{equation}
    \label{eq:D-both}
    \mathcal{D}_\mu\Phi^I{}_A
    =\partial_\mu\Phi^I{}_A
     -i(A_\mu)^I{}_J\Phi^J{}_A
     +\omega_\mu{}^B{}_A\Phi^I{}_B.
  \end{equation}
  The connection \(\omega_\mu{}^{AB}\) is, at this stage, an
  independent gauge field.
\end{assumption}

Assumption~\ref{ass:local-lorentz} is the step that makes \(A\) a
genuine tangent index rather than an internal flavour that happens
to have four values. Combined with Assumption~\ref{ass:lorentzian},
it is the origin of local Lorentz invariance in NAMCG. It is not
derived from Presentation~A.

\section{The vacuum and the broken generators}
\label{sec:ch10-broken}

The polar vacuum of Chapter~4, in the Minkowski frame,
\begin{equation}
  \label{eq:vac-lock}
  \langle\Phi^I{}_A\rangle=v\,\delta^I{}_A,
\end{equation}
is invariant only under the diagonal combination
\begin{equation}
  \label{eq:diagonal}
  \alpha^I{}_J
  =-\,v^{-1}\delta^I{}_B\,\beta^B{}_C\,\delta^C{}_J
   \Big|_{\text{locked \(4\times 4\) block}},
\end{equation}
that is, under simultaneous internal and Lorentz rotations that
rotate the locked axes together. If \(G_{\mathrm{int}}\) contains
\(\mathrm{SO}(3,1)\) (or a group large enough to have an
\(\mathfrak{so}(3,1)\) subalgebra acting on the locked indices),
six generators are unbroken and six are broken.

\begin{proposition}[Goldstone count]
  \label{prop:goldstone-count}
  About the vacuum~\eqref{eq:vac-lock}, with
  Assumption~\ref{ass:local-lorentz} in force and with the locked
  block of dimension four, there are six broken generators. The
  corresponding Goldstone fields are the antisymmetric components
  \(\omega_{AB}\) of~\eqref{eq:linear-frame}.
\end{proposition}

\begin{proof}
  The vacuum stabilizer is the diagonal
  \(\mathrm{SO}(3,1)_{\mathrm{diag}}\). The coset
  \(\mathrm{SO}(3,1)_{\mathrm{Lorentz}}\times
   \mathrm{SO}(3,1)_{\mathrm{int}}/\mathrm{SO}(3,1)_{\mathrm{diag}}\)
  has dimension \(6+6-6=6\). An infinitesimal broken transformation
  with parameter \(\beta_{AB}\) shifts
  \(\Phi^I{}_A\) by \(v\beta^I{}_A\), which is antisymmetric in
  \((I,A)\) after lowering with \(\eta\). That shift is the
  fluctuation \(\omega_{AB}\).
\end{proof}

If \(G_{\mathrm{int}}\) is larger than \(\mathrm{SO}(3,1)\),
additional Goldstones live in the unlocked block. They belong to
Assumption~\ref{ass:locked-block} and are taken heavy.

\section{Eating \(\omega_{AB}\)}
\label{sec:ch10-eat}

A local broken symmetry with an independent connection is a Higgs
mechanism. Introduce the Stueckelberg field
\(\Lambda(x)\in\mathrm{SO}(3,1)\) already present in the polar
decomposition~\eqref{eq:polar}, and replace the spin connection by
the dressed connection
\begin{equation}
  \label{eq:dressed-omega}
  \omega_\mu
  \;\longmapsto\;
  \Lambda^{-1}\omega_\mu\Lambda+\Lambda^{-1}\partial_\mu\Lambda.
\end{equation}
Choose the unitary gauge
\begin{equation}
  \label{eq:unitary}
  \Lambda(x)=\mathbb{I},
\end{equation}
which is admissible wherever the polar decomposition is smooth
(that is, wherever Assumption~\ref{ass:nondeg} holds and no
condensate defect forces \(\det e=0\)). In that gauge
\(\omega_{AB}\) has been absorbed into \(\omega_\mu{}^{AB}\), and
the kinetic term of \(\Phi\) produces a mass term
\begin{equation}
  \label{eq:omega-mass}
  \mathcal{L}_{\omega^2}
  \supset
  v^2\,\omega_\mu{}^{AB}\omega^\mu{}_{AB}.
\end{equation}
The connection is massive of order \(v\), or it is identified with
the Levi-Civita spin connection of the composite frame and the
mass term is reinterpreted as the statement that torsion is
gapped.

\begin{theorem}[Lorentz locking]
  \label{thm:locking}
  Under Assumptions~\ref{ass:frame-lock},
  \ref{ass:local-lorentz}, and~\ref{ass:nondeg}, there exists a
  gauge in which
  \begin{enumerate}
    \item the six Goldstones \(\omega_{AB}\) do not propagate;
    \item the composite object \(e^A_\mu\) of
          Definition~\ref{def:vierbein} transforms only under
          the diagonal local Lorentz group, as a vierbein;
    \item the independent spin connection is either massive of
          order \(v\) or slaved to \(e^A_\mu\) by the torsion-free
          condition~\eqref{eq:spin-connection}.
  \end{enumerate}
\end{theorem}

\begin{proof}
  Proposition~\ref{prop:goldstone-count} supplies six Goldstones.
  The Stueckelberg replacement~\eqref{eq:dressed-omega} and the
  gauge~\eqref{eq:unitary} remove them from the spectrum. In that
  gauge \(\Phi^{IA}=(v+\sigma)\delta^{IA}\) plus tensor and vector
  blocks, so \(e^A_\mu\) defined from \(\partial_\mu\Phi\) or from
  the normalized vacuum transforms as
  \(e^A_\mu\mapsto\Lambda_{\mathrm{diag}}{}^A{}_B e^B_\mu\). The
  mass term~\eqref{eq:omega-mass} gaps any component of
  \(\omega_\mu{}^{AB}\) that is not slaved to \(e\). The slaved
  remainder is uniquely the torsion-free spin connection of
  \(e^A_\mu\), by the standard Cartan lemma used
  in~\eqref{eq:spin-connection}.
\end{proof}

\begin{remark}[Defects]
  Unitary gauge fails where \(\Lambda(x)\) cannot be set to the
  identity continuously---at defects where \(\det e=0\). Those loci
  are condensate defects, not coordinate singularities of an
  already existing manifold. They are possible ultraviolet
  endpoints of the regular cores of Chapter~13, and they are not
  analyzed further here.
\end{remark}

\section{The internal vector after locking}
\label{sec:ch10-vector}

The internal connection \(A_\mu\) of~\eqref{eq:Dmu-Phi} is not the
same object as \(\omega_\mu{}^{AB}\). After locking, four of its
components in the locked block mix with \(\omega_\mu\). The
orthogonal combination is an ordinary Yang--Mills field of the
unbroken diagonal subgroup, or of whatever remains of
\(G_{\mathrm{int}}\) outside \(\mathrm{SO}(3,1)\).

Its mass is a standard Higgs mass,
\begin{equation}
  \label{eq:A-mass}
  m_A\sim g_{\mathrm{YM}}v,
\end{equation}
unless a residual exact gauge symmetry is left unbroken. In the
latter case the infrared theory contains general relativity plus
Yang--Mills, which is phenomenologically welcome and does not
disturb Theorem~\ref{thm:mT}. In the former case the vector block
decouples at scale \(m_A\) and is irrelevant below \(v\).

This closes the vector line of the counting table.

\section{Diffeomorphisms}
\label{sec:ch10-diffs}

Diffeomorphisms act on the coordinate index of every tensor
constructed from \(\Phi(x)\). They are not eaten. They are the
gauge symmetry of the composite metric, and they are what force
\(\nabla^\mu T_{\mu\nu}^{\mathrm{total}}=0\) in Chapter~11.

After Theorem~\ref{thm:locking}, an infinitesimal diffeomorphism
\(\xi^\mu\) acts on the vierbein in the standard hybrid way,
\begin{equation}
  \label{eq:lie-vierbein}
  \delta_\xi e^A_\mu
  =\xi^\rho\partial_\rho e^A_\mu
   +\partial_\mu\xi^\rho\,e^A_\rho,
\end{equation}
up to a local Lorentz transformation that can be used to keep a
chosen gauge on the frame. The induced action on \(g_{\mu\nu}\) is
the Lie derivative~\eqref{eq:lie-g}. No extra scalar Goldstone of
diffeomorphisms appears: diffeomorphisms are gauged, not broken, in
the condensed phase.

\section{Completed counting table}
\label{sec:ch10-table}

\begin{center}
\begin{tabular}{@{}llll@{}}
\toprule
Stage & What is removed & Remaining & Where \\
\midrule
algebraic split & --- & 16 & Ch.~5 \\
eat \(\omega_{AB}\) & 6 Lorentz Goldstones & 10 & this chapter \\
gap or Higgs \(A_\mu\) & 4 vector & 6 & this chapter \\
gap \(\psi\) & 1 conformal scalar & 5 & Assp.~\ref{ass:gap-psi} \\
\(m_T=0\) and TT gauge & 3 residual constraints & 2 TT & Ch.~7 \\
\bottomrule
\end{tabular}
\end{center}

The six remaining after the vector is removed are the ten metric
components minus four diffeomorphisms. Gapping \(\psi\) and imposing
massless TT gauge then yields two polarizations. The table of
Section~\ref{sec:ch5-count} is thereby completed, with the one
outstanding dynamical hypothesis still being
Assumption~\ref{ass:gap-psi}.

\section{What locking does not do}
\label{sec:ch10-not}

Theorem~\ref{thm:locking} does not derive Lorentzian signature
(Assumption~\ref{ass:lorentzian} remains). It does not derive local
Lorentz invariance from Presentation~A
(Assumption~\ref{ass:local-lorentz} remains). It does not gap
\(\psi\). It does not compute \(\Geff\). It does not conserve
\(\Tent_{\mu\nu}\).

What it does is make the later covariant derivatives honest. From
here on, \(\nabla_\mu\) is the Levi-Civita connection of the
composite metric, and statements such as
\(\nabla^\mu T_{\mu\nu}^{\mathrm{total}}=0\) are tensor equations
on the emergent spacetime rather than slogans.

\section{Exercises}
\label{sec:ch10-exercises}

\begin{enumerate}
  \item \textbf{Stabilizer.}
    Verify that~\eqref{eq:vac-lock} is invariant under the
    diagonal combination~\eqref{eq:diagonal} and under no other
    independent pair \((\alpha,\beta)\) acting on the locked block.
  \item \textbf{Coset dimension.}
    Compute
    \(\dim\bigl(\mathrm{SO}(3,1)\times\mathrm{SO}(3,1)
    /\mathrm{SO}(3,1)\bigr)\) and match it to the six components
    of \(\omega_{AB}\).
  \item \textbf{Stueckelberg.}
    Starting from~\eqref{eq:D-both} and~\eqref{eq:polar}, perform
    the replacement~\eqref{eq:dressed-omega} and isolate the
    quadratic term~\eqref{eq:omega-mass} in unitary gauge.
  \item \textbf{Torsion-free slaving.}
    Given \(e^A_\mu\), write the explicit formula for the
    torsion-free spin connection
    \(\omega_\mu{}^{AB}(e)\). Confirm that it transforms
    inhomogeneously under local Lorentz transformations, as a
    connection must.
  \item \textbf{Defect.}
    Construct a hedgehog configuration of \(\Lambda(x)\) on a
    spatial 2-sphere inside which \(\det e=0\) at a point. Why
    can unitary gauge not be imposed globally on that slice?
  \item \textbf{Residual Yang--Mills.}
    Suppose \(G_{\mathrm{int}}=\mathrm{SO}(3,1)\times H\) with
    \(H\) unbroken. What is the infrared field content after
    Theorem~\ref{thm:locking}, in addition to the graviton?
  \item \textbf{Hybrid diffeomorphism.}
    Derive~\eqref{eq:lie-vierbein} from the action of a
    coordinate shift on \(\Phi(x)\) and
    Definition~\ref{def:vierbein}.
\end{enumerate}

\section{Bridge to Chapter~11}
\label{sec:ch10-bridge}

The emergent geometry now has a metric, a torsion-free connection,
two massless graviton polarizations, and a computed \(\Geff\).
Diffeomorphism invariance is in force. Noether's second theorem
therefore requires
\begin{equation}
  \label{eq:cons-next}
  \nabla^\mu T_{\mu\nu}^{\mathrm{total}}=0
\end{equation}
of whatever stress tensor sources the Einstein equation. Chapter~11
defines \(T_{\mu\nu}^{\mathrm{total}}\) as the sum of ordinary
matter, the tensor-block stress already used in Chapter~7, and the
entropic stress \(\Tent_{\mu\nu}\) obtained by varying
\(U_{\mathrm{ent}}\). It then proves that the sum is divergenceless
on shell. This conservation statement constrains the total stress
when it is derived from the specified diffeomorphism-invariant
action. It does not, by itself, derive the phenomenological drag law
or show that a reduced drag model is its covariant limit; that
connection remains Problem~\ref{sec:status-I}.